\section*{Homework for Section 7.2: 40 Points}
\addcontentsline{toc}{section}{Homework for Section 7.2: 40 Points}
\vspace{-0.5cm}
In Problems 1 - 30, use appropirate algebra and \hyperlink{def:Theorem 7.2.1}{\textbf{Theorem 7.2.1}}
to find the given inverse Laplace transform $\mathscr{L}^{-1} \{ f(t) \}$.

\phantomsection
  \addcontentsline{toc}{subsection}{\textbf{Theorem 7.2.1}}
  \begin{boxedtheorem}[7.2.1] \hypertarget{def:Theorem 7.2.1}{}
  \ \newline
    \textbf{Some Inverse Transforms }
    \begin{enumerate}[label=(\alph*)]
      \item $\displaystyle 1 = \mathscr{L}^{-1} \left\{ \frac{1}{s} \right\}$
      \item $\displaystyle t^n = \mathscr{L}^{-1} \left\{ \frac{n!}{s^{n+1}} \right\}, \quad n = 1, 2, 3, \dots$
      \item $\displaystyle e^{at} = \mathscr{L}^{-1} \left\{ \frac{1}{s - a} \right\}$
      \item $\displaystyle \sin kt = \mathscr{L}^{-1} \left\{ \frac{k}{s^2 + k^2} \right\}$
      \item $\displaystyle \cos kt = \mathscr{L}^{-1} \left\{ \frac{s}{s^2 + k^2} \right\}$
      \item $\displaystyle \sinh kt = \mathscr{L}^{-1} \left\{ \frac{k}{s^2 - k^2} \right\}$
      \item $\displaystyle \cosh kt = \mathscr{L}^{-1} \left\{ \frac{s}{s^2 - k^2} \right\}$
    \end{enumerate}
  \end{boxedtheorem}
  

\subsection*{Section 7.2: 16 Points}
\addcontentsline{toc}{subsection}{Section 7.2: 16 Points}

  \subsubsection*{Problem 7.2.4}
  \addcontentsline{toc}{subsubsection}{Problem 7.2.4}
  \vspace{-0.5cm}

  \begin{flalign*}
      & \text{Given: } \mathscr{L}^{-1} \left\{ \left(\frac{2}{s}-\frac{1}{s^{3}}\right)^{2} \right\}, \ f(t) = \left(\frac{2}{s}-\frac{1}{s^{3}}\right)^{2} & \\
      & \text{Find:  } \mathscr{L}^{-1} \left\{ f(t) \right\} &
  \end{flalign*}
  \begin{solution}
  Using the table from \hyperlink{def:Theorem 7.2.1}{\textbf{Theorem 7.2.1}}, we can use 
  the linear property of the inverse Laplacian and evaluate. Here we will expand out the power
  and use power rule.

  \begin{align*}
    \mathscr{L}^{-1} \{ f(t) \} &= \mathscr{L}^{-1} \left\{ \left(\frac{2}{s}-\frac{1}{s^{3}}\right)^{2} \right\} \\
                                &= \mathscr{L}^{-1} \left\{ \frac{4}{s^{2}}-\frac{4}{s^{4}}+\frac{1}{s^{6}} \right\} \\
                                &= 4\mathscr{L}^{-1} \left\{ \frac{1}{s^{2}}\right\} -4 \mathscr{L}^{-1} \left\{\frac{1}{s^{4}} \right\} + \mathscr{L}^{-1} \left\{\frac{1}{s^{6}} \right\} \\
                                &= \frac{4}{1!}\mathscr{L}^{-1} \left\{ \frac{1!}{s^{1+1}}\right\} - \frac{4}{3!} \mathscr{L}^{-1} \left\{\frac{3!}{s^{3+1}} \right\} + \frac{1}{5!} \mathscr{L}^{-1} \left\{\frac{5!}{s^{5+1}} \right\} \\
                                &= 4t-\frac{4}{6}t^{3}+\frac{1}{120}t^{5}
  \end{align*}
  % \mathscr{L}^{-1} \{ f(t) \} &= 4t-\frac{2}{3}t^{3}+\frac{1}{120}t^{5}
  \begin{equation*}
    \Longrightarrow \boxed{\mathscr{L}^{-1}\{f(t)\} = 4t-\frac{2}{3}t^{3}+\frac{1}{120}t^{5}, \ t \geq 0}
  \end{equation*}

  \end{solution}

  \subsubsection*{Problem 7.2.8}
  \addcontentsline{toc}{subsubsection}{Problem 7.2.8}
  \vspace{-0.5cm}

  \begin{flalign*}
      & \text{Given: } \mathscr{L}^{-1} \left\{ \frac{4}{s}+\frac{6}{s^{5}}-\frac{1}{s+8} \right\}, \ f(t) = \frac{4}{s}+\frac{6}{s^{5}}-\frac{1}{s+8} & \\
      & \text{Find:  } \mathscr{L}^{-1} \left\{ f(t) \right\} &
  \end{flalign*}
  \begin{solution}
  Using the table from \hyperlink{def:Theorem 7.2.1}{\textbf{Theorem 7.2.1}}, we can use 
  the linear property of the inverse Laplacian and evaluate. Here we will expand out the power
  and use power rule.

  \begin{align*}
    \mathscr{L}^{-1} \{ f(t) \} &= \mathscr{L}^{-1} \left\{ \frac{4}{s}+\frac{6}{s^{5}}-\frac{1}{s+8} \right\} \\
                                &= 4\mathscr{L}^{-1} \left\{ \frac{1}{s}\right\} + 6\mathscr{L}^{-1} \left\{\frac{1}{s^{5}} \right\} - \mathscr{L}^{-1} \left\{\frac{1}{s+8} \right\} \\
                                &= 4\mathscr{L}^{-1} \left\{ \frac{1}{s}\right\} +  \frac{6}{4!} \mathscr{L}^{-1} \left\{\frac{4!}{s^{4+1}} \right\} - \mathscr{L}^{-1} \left\{\frac{1}{s+8} \right\} \\
    \mathscr{L}^{-1} \{ f(t) \} &= 4 + \frac{6}{24}t^{4}-e^{-8t}
  \end{align*}
  %     \mathscr{L}^{-1} \{ f(t) \} &= 4 + \frac{1}{4}t^{4}-e^{-8t}

  \begin{equation*}
    \Longrightarrow \boxed{\mathscr{L}^{-1}\{f(t)\} = 4 + \frac{1}{4}t^{4}-e^{-8t}, \ t \geq 0}
  \end{equation*}

  \end{solution}

  \subsubsection*{Problem 7.2.16}
  \addcontentsline{toc}{subsubsection}{Problem 7.2.16}
  \vspace{-0.5cm}

  \begin{flalign*}
      & \text{Given: } \mathscr{L}^{-1} \left\{ \frac{s+1}{s^{2}+2} \right\}, \ f(t) = \frac{s+1}{s^{2}+2} & \\
      & \text{Find:  } \mathscr{L}^{-1} \left\{ f(t) \right\} &
  \end{flalign*}
  \begin{solution}
  Using the table from \hyperlink{def:Theorem 7.2.1}{\textbf{Theorem 7.2.1}}, we can use 
  the linear property of the inverse Laplacian and evaluate. Here we will expand out the power
  and use power rule.

  \begin{align*}
    \mathscr{L}^{-1} \{ f(t) \} &= \mathscr{L}^{-1} \left\{ \frac{s+1}{s^{2}+2} \right\} \\
                                &= \mathscr{L}^{-1} \left\{ \frac{s}{s^{2}+2} + \frac{1}{s^{2}+2}\right\} \\
                                &= \mathscr{L}^{-1} \left\{ \frac{s}{s^{2}+\left(\sqrt{2}\right)^2} \right\} + \frac{1}{\sqrt{2}} \mathscr{L}^{-1} \left\{ \left(\frac{\sqrt{2}}{s^{2}+\left(\sqrt{2}\right)^2}\right) \right\} \\
                                &= \cos\left(t\sqrt{2}\right) + \frac{1}{\sqrt{2}} \sin\left(t\sqrt{2}\right) \\
    \mathscr{L}^{-1} \{ f(t) \} &= \cos\left(t\sqrt{2}\right) + \frac{\sqrt{2}}{2} \sin\left(t\sqrt{2}\right)
  \end{align*}

  \begin{equation*}
    \Longrightarrow \boxed{\mathscr{L}^{-1} \{ f(t) \} = \cos\left(t\sqrt{2}\right) + \frac{\sqrt{2}}{2} \sin\left(t\sqrt{2}\right), \ t \geq 0}
  \end{equation*}

  \end{solution}

  \subsubsection*{Problem 7.2.24}
  \addcontentsline{toc}{subsubsection}{Problem 7.2.24}
  \vspace{-0.5cm}

  \begin{flalign*}
      & \text{Given: } \mathscr{L}^{-1} \left\{ \frac{s^{2}+1}{s\left(s-1\right)\left(s+1\right)\left(s-2\right)} \right\}, \ f(t) = \frac{s^{2}+1}{s\left(s-1\right)\left(s+1\right)\left(s-2\right)} & \\
      & \text{Find:  } \mathscr{L}^{-1} \left\{ f(t) \right\} &
  \end{flalign*}
  \begin{solution}
  Using the table from \hyperlink{def:Theorem 7.2.1}{\textbf{Theorem 7.2.1}}, we can use 
  the linear property of the inverse Laplacian and evaluate. Here we will expand out the power
  and use power rule.

  \begin{align*}
    \mathscr{L}^{-1} \{ f(t) \} &= \mathscr{L}^{-1} \left\{ \frac{s^{2}+1}{s\left(s-1\right)\left(s+1\right)\left(s-2\right)} \right\} \\
                                &= \mathscr{L}^{-1} \left\{\frac{A}{s}+\frac{B}{s-1}+\frac{C}{s+1}+\frac{D}{s-2} \right\}
  \end{align*}
  Using Heaviside cover up method.
  \begin{align*}
    & A=\frac{1}{2}, \ B=-1, \ C=-\frac{1}{3}, \ D=\frac{5}{6} \\
                \Longrightarrow &= \mathscr{L}^{-1} \left\{ \frac{\frac{1}{2}}{s}+\frac{-1}{s-1}+\frac{-\frac{1}{3}}{s+1}+\frac{\frac{5}{6}}{s-2} \right\} \\
                                &= \frac{1}{2} - \mathscr{L}^{-1} \left\{\frac{1}{s-1} \right\} -\frac{1}{3} \mathscr{L}^{-1} \left\{\frac{1}{s+1} \right\} +\frac{5}{6} \mathscr{L}^{-1} \left\{\frac{1}{s-2} \right\}\\
    \mathscr{L}^{-1} \{ f(t) \} &= \frac{1}{2}-e^{t}-\frac{1}{3}e^{-t}+\frac{5}{6}e^{2t}
  \end{align*}

  \begin{equation*}
    \Longrightarrow \boxed{\mathscr{L}^{-1} \{ f(t) \} = \frac{1}{2}-e^{t}-\frac{1}{3}e^{-t}+\frac{5}{6}e^{2t}, \ t \geq 0}
  \end{equation*}

  \end{solution}

\subsection*{Section 7.2: 24 Points}
\addcontentsline{toc}{subsection}{Section 7.2: 24 Points}
In Problems 35 - 44, use the Laplace transform to solve the given initial value problem.

It should be noted the formula for $\mathscr{L} \{ f^{\left(n\right)}\left(t\right) \}$.

\begin{boxedtheorem}[7.2.2] \hypertarget{def:Theorem 7.2.2}{}
  \ \newline
  \textbf{Transform of a Derivative} \newline
  If $f, f', \dots, f^{(n-1)}$ are continuous on $[0, \infty)$ and are of exponential order
  and if $f^{(n)}(t)$ is piecewise continuous on $[0, \infty)$, then

  \vspace{-0.5cm}
  \begin{align}
    \mathscr{L}\{f^{(n)}(t)\} &= s^n F(s) - s^{n-1} f(0) - s^{n-2} f'(0) - \dots - f^{(n-1)}(0) \\
                              &= s^n F(s) -\sum_{k=1}^{n}s^{\left(k-1\right)}f^{\left(n-k\right)}(0),
  \end{align}

  where $F(s) = \mathscr{L}\{f(t)\}$. 
\end{boxedtheorem}
\addcontentsline{toc}{subsection}{\textbf{Theorem 7.2.2}}

  \newpage
  \subsubsection*{Problem 7.2.36}
  \addcontentsline{toc}{subsubsection}{Problem 7.2.36}
  \vspace{-0.5cm}

  \begin{flalign*}
      & \text{Given: } 2\frac{dy}{dt}+y=0, \ y(0)=-3 & \\
      & \text{Find:  } y(t) &
  \end{flalign*}
  \begin{solution}
  The main process to apply the Laplace transform, using \hyperlink{def:Theorem 7.1.1}{\textbf{Theorem 7.1.1}} and \hyperlink{def:Theorem 7.2.2}{\textbf{Theorem 7.2.2}}, over the entire equation,
  isolate $Y(s)$, and apply the inverse Laplace transform using \hyperlink{def:Theorem 7.2.1}{\textbf{Theorem 7.2.1}}.

  \begin{align*}
                                    & 2\frac{dy}{dt}+y  = 0 \\
    \mathscr{L} \hspace{-0.10cm}: \ & 2\frac{dy}{dt}+y  = 0 \\
                    \Longrightarrow & 2\left(-y(0)+sY(s)\right)+Y(s) = 0 
  \end{align*}
  Rearrange and isolate $Y(s)$
  \begin{align*}
    \left(2s+1\right) Y(s) &= 2y(0) \\
                      Y(s) &= \frac{2y(0)}{\left(2s+1\right)} \\
                      Y(s) &= \frac{2\left(-3\right)}{\left(2s+1\right)} \\
                      Y(s) &= -\frac{6}{\left(2s+1\right)} \\
                      Y(s) &= -\frac{6}{2\left(s+\frac{1}{2}\right)} \\
                      Y(s) &= -\frac{3}{\left(s+\frac{1}{2}\right)}
  \end{align*}
  Apply the inverse Laplace transform.
  \begin{align*}
     \mathscr{L}^{-1} \hspace{-0.10cm}: \ & Y(s) = -\frac{3}{\left(s+\frac{1}{2}\right)} \\
                                        \Longrightarrow  & \boxed{y(t) = -3e^{-\frac{1}{2}t}}
  \end{align*}

  \end{solution}

  \newpage
  \subsubsection*{Problem 7.2.40}
  \addcontentsline{toc}{subsubsection}{Problem 7.2.40}
  \vspace{-0.5cm}

  \begin{flalign*}
      & \text{Given: } y''-4y'=6e^{3t}-3e^{-t}, \ y(0) = 1, \ y'(0) = -1 & \\
      & \text{Find:  } y(t) &
  \end{flalign*}
  \begin{solution}
  The main process to apply the Laplace transform, using \hyperlink{def:Theorem 7.1.1}{\textbf{Theorem 7.1.1}} and \hyperlink{def:Theorem 7.2.2}{\textbf{Theorem 7.2.2}}, over the entire equation,
  isolate $Y(s)$, and apply the inverse Laplace transform using \hyperlink{def:Theorem 7.2.1}{\textbf{Theorem 7.2.1}}.

  \begin{align*}
                                    & y''-4y'=6e^{3t}-3e^{-t} \\
    \mathscr{L} \hspace{-0.10cm}: \ & y''-4y'=6e^{3t}-3e^{-t} \\
                    \Longrightarrow & \left(-y'(0)-sy(0)+s^{2}Y(s)\right)-4\left(-y(0)+sY(s)\right)=\frac{6}{s-3}-\frac{3}{s+1}
  \end{align*}
  Rearrange and isolate $Y(s)$.
  \begin{align*}
          -y'(0)-sy(0)+4y(0)-4s+s^{2}Y(s)-4sY(s) &= \frac{6}{s-3}-\frac{3}{s+1} \\
         \left(s^{2}-4s\right) Y(s) &= \frac{6}{s-3}-\frac{3}{s+1}+y'(0)+y(0)\left(s-4\right) \\
         \left(s^{2}-4s\right) Y(s) &= \frac{6}{s-3}-\frac{3}{s+1}+\left(-1\right)+\left(1\right)\left(s-4\right) \\
          \left(\left(s\right)\left(s-4\right)\right) Y(s) &= \frac{6}{s-3}-\frac{3}{s+1}+\left(s-5\right) \\
          Y(s) &= \frac{6}{s\left(s-3\right)\left(s-4\right)}-\frac{3}{s\left(s+1\right)\left(s-4\right)}+\frac{s-5}{s\left(s-4\right)} \\
          Y(s) &= \frac{A}{s}+\frac{B}{s-3}+\frac{C}{s-4}+\frac{D}{s+1}
  \end{align*}
  Using Heaviside cover up method.
  \begin{align*}
    & A=\frac{5}{2}, \ B=-2, \ C=\frac{11}{10}, \ D=-\frac{3}{5} \\
                \Longrightarrow Y(s) &= \frac{\frac{5}{2}}{s}+\frac{-2}{s-3}+\frac{\frac{11}{10}}{s-4}+\frac{-\frac{3}{5}}{s+1}
  \end{align*}
  Apply the inverse Laplace transform.
  \begin{align*}
    \mathscr{L}^{-1} \hspace{-0.10cm}: \ & Y(s) = \frac{\frac{5}{2}}{s}+\frac{-2}{s-3}+\frac{\frac{11}{10}}{s-4}+\frac{-\frac{3}{5}}{s+1} \\
                                         & y(t) = \frac{5}{2} \mathscr{L}^{-1} \left\{\frac{1}{s} \right\} - 2\mathscr{L}^{-1} \left\{\frac{1}{s-3} \right\} + \frac{11}{10} \mathscr{L}^{-1} \left\{\frac{1}{s-4} \right\} - \frac{3}{5} \mathscr{L}^{-1} \left\{\frac{1}{s+1} \right\} \\
                       \Longrightarrow   & \boxed{y(t) = \frac{5}{2}-2e^{3t}+\frac{11}{10}e^{4t}-\frac{3}{5}e^{-t}}
  \end{align*}

  \end{solution}

\newpage
In Problems 45 and \textbf{46} use the Laplace transform and these inverses to solve the given initial-value problem.

  \subsubsection*{Problem 7.2.46}
  \addcontentsline{toc}{subsubsection}{Problem 7.2.46}
  \vspace{-0.5cm}

  \begin{flalign*}
      & \text{Given: } y''-2y'+5y=0, \ y(0) = 1, \ y'(0) = 3 & \\
      & \text{Find:  } y(t) &
  \end{flalign*}
  \begin{solution}
  The main process to apply the Laplace transform, using \hyperlink{def:Theorem 7.1.1}{\textbf{Theorem 7.1.1}} and \hyperlink{def:Theorem 7.2.2}{\textbf{Theorem 7.2.2}}, over the entire equation,
  isolate $Y(s)$, and apply the inverse Laplace transform using \hyperlink{def:Theorem 7.2.1}{\textbf{Theorem 7.2.1}}.

  \begin{align*}
                                    & y''-2y'+5y=0 \\
    \mathscr{L} \hspace{-0.10cm}: \ & y''-2y'+5y=0 \\
                    \Longrightarrow & \left(-y'(0)-sy(0)+s^{2}Y(s)\right)-2\left(-y(0)+sY(s)\right) + 5Y(s)=0
  \end{align*}
  Rearrange and isolate $Y(s)$.
  \begin{align*}
          \left(s^{2}-2s+5\right)Y(s) &= y'(0)+y(0)\left(s-2\right) \\
          \left(s^{2}-2s+5\right)Y(s) &= \left(3\right)+\left(1\right)\left(s-2\right) \\
          \left(s^{2}-2s+5\right)Y(s) &= s+1 \\
          Y(s) &= \frac{s+1}{s^{2}-2s+5}
  \end{align*}
  Completing the square.
  \begin{align*}
    Y(s) &= \frac{s+1}{\left(s-1\right)^{2}+4} \\
    Y(s) &= \frac{s-1}{\left(s-1\right)^{2}+4}+\frac{2}{\left(s-1\right)^{2}+4}
  \end{align*}
  Apply the inverse Laplace transform alongside the translation \hyperlink{def:Theorem 7.3.1}{\textbf{Theorem 7.3.1}}.
  \begin{align*}
    \mathscr{L}^{-1} \hspace{-0.10cm}: \ & Y(s) = \frac{s-1}{\left(s-1\right)^{2}+2^{2}}+\frac{2}{\left(s-1\right)^{2}+2^{2}} \\
                                         & y(t) = e^{t}\cos\left(2t\right)+e^{t}\sin\left(2t\right) \\
                       \Longrightarrow   & \boxed{y(t) = e^{t}\left(\cos\left(2t\right)+\sin\left(2t\right)\right)}
  \end{align*}

  \end{solution}
